\documentclass[11pt,a4paper]{article}

\usepackage{fontspec}
\setmainfont{Times New Roman}
\setsansfont{Arial}
\setmonofont{Menlo}

\usepackage[a4paper,margin=24mm]{geometry}
\usepackage{amsmath,amssymb,mathtools}
\usepackage{booktabs}
\usepackage{longtable}
\usepackage{enumitem}
\usepackage{xcolor}
\usepackage{hyperref}
\usepackage{xurl}
\hypersetup{
  colorlinks=true,
  linkcolor=black,
  urlcolor=blue,
  citecolor=black
}

\setlist[itemize]{leftmargin=*,topsep=3pt,itemsep=2pt}
\setlist[enumerate]{leftmargin=*,topsep=3pt,itemsep=2pt}
\renewcommand{\arraystretch}{1.18}
\emergencystretch=3em

\title{\textbf{Формальное описание работы Particle Life Plus}}
\author{Описание по реализации в репозитории ASAL}
\date{11 апреля 2026}

\begin{document}
\maketitle

\section*{Область действия документа}

Документ формализует субстрат \texttt{ParticleLifePlus}.
Основная реализация находится в \path{substrates/plife_plus.py}.
Для сравнения используется классический \texttt{ParticleLife} из
\path{substrates/plife.py}. Параметры запуска \texttt{plife\_plus} в фабрике
субстратов задаются в \path{substrates/__init__.py}, а CLI-параметры перечислены
в \path{scripts/util.py}.

Particle Life Plus является системой из $N$ частиц в двумерной области. В отличие
от классического Particle Life, где цвет частицы является дискретным классом,
здесь цвет является непрерывным вектором. Пара взаимодействующих цветовых
векторов подается в небольшую нейросеть, которая одновременно задает амплитуду
механического взаимодействия и локальное изменение цвета.

\section*{1. Обозначения}

Пусть:
\begin{itemize}
  \item $N$ --- число частиц, \texttt{n\_particles};
  \item $K$ --- размерность цветового пространства, \texttt{n\_colors};
  \item $d=2$ --- размерность физического пространства;
  \item $L$ --- размер стороны квадратной области, \texttt{world\_size};
  \item $m>0$ --- масса частицы, \texttt{mass};
  \item $\Delta t>0$ --- шаг интегрирования, \texttt{dt};
  \item $h>0$ --- период полураспада скорости, \texttt{half\_life};
  \item $r_{\max}>0$ --- радиус взаимодействия, \texttt{rmax};
  \item $\beta\in(0,1)$ --- граница внутренней зоны отталкивания, \texttt{beta};
  \item $\theta$ --- параметры нейросети \texttt{PLifeNetwork}.
\end{itemize}

Состояние системы в дискретный момент $n$:
\[
S_n=\{(x_i^n,v_i^n,c_i^n)\}_{i=1}^{N},
\]
где
\[
x_i^n\in[0,L]^2,\qquad v_i^n\in\mathbb{R}^2,\qquad
c_i^n\in\mathbb{R}^K,\qquad \|c_i^n\|_2=1.
\]

\section*{2. Инициализация}

Для каждой частицы независимо генерируются:
\[
\xi_i \sim \mathcal{N}(0,I_K),\qquad
c_i^0=\frac{\xi_i}{\|\xi_i\|_2},
\]
\[
x_i^0 \sim \mathrm{Uniform}([0,L]^2),\qquad
v_i^0=0.
\]

Параметры $\theta$ нейросети создаются стандартной инициализацией Flax Linen:
\[
\theta = \mathrm{init}_{\mathrm{Flax}}(\mathrm{PLifeNetwork}).
\]

\section*{3. Нейросеть взаимодействия}

Для упорядоченной пары частиц $(i,j)$ нейросеть получает конкатенацию двух
цветовых векторов:
\[
z_{ij} = [c_i,c_j]\in\mathbb{R}^{2K}.
\]

Архитектура \texttt{PLifeNetwork}:
\[
\mathbb{R}^{2K}
\xrightarrow{\mathrm{Dense}(8),\,\tanh}
\mathbb{R}^{8}
\xrightarrow{\mathrm{Dense}(8),\,\tanh}
\mathbb{R}^{8}
\xrightarrow{\mathrm{Dense}(8),\,\tanh}
\mathbb{R}^{8}
\xrightarrow{\mathrm{Dense}(1+K)}
\mathbb{R}^{1+K}.
\]

Пусть последний слой выдает вектор
\[
y_{ij}=(y_{ij}^{(0)},y_{ij}^{(1:K)}).
\]
Тогда используются две величины:
\[
\alpha_{ij}=1.5\,\tanh(y_{ij}^{(0)})\in[-1.5,1.5],
\]
\[
g_{ij}=\tanh(y_{ij}^{(1:K)})\in[-1,1]^K.
\]

Здесь $\alpha_{ij}$ управляет амплитудой силы в средней зоне взаимодействия, а
$g_{ij}$ задает направление изменения цвета частицы $i$ под влиянием частицы $j$.

\section*{4. Геометрия пары частиц}

Для пары $(i,j)$ сначала вычисляется вектор смещения:
\[
\delta_{ij}=x_j-x_i.
\]

Если граница \texttt{border='torus'}, применяется правило минимального образа
покомпонентно:
\[
\delta_{ij,k}=
\begin{cases}
\delta_{ij,k}-L, & \delta_{ij,k}>L/2,\\
\delta_{ij,k}+L, & \delta_{ij,k}<-L/2,\\
\delta_{ij,k}, & \text{иначе}.
\end{cases}
\]

Если граница \texttt{border='wall'}, при вычислении парной силы используется
обычное смещение $x_j-x_i$ без периодической поправки.

Далее:
\[
\rho_{ij}=\|\delta_{ij}\|_2,\qquad
e_{ij}=\frac{\delta_{ij}}{\rho_{ij}+\varepsilon},
\]
где в коде $\varepsilon=10^{-8}$.

\section*{5. Граф силы}

Вводится безразмерная дистанция
\[
q_{ij}=\frac{\rho_{ij}}{r_{\max}}.
\]

Скалярный граф силы:
\[
\Phi(q;\alpha,\beta)=
\begin{cases}
\dfrac{q}{\beta}-1, & q<\beta,\\[8pt]
\alpha\left(1-\dfrac{|2q-1-\beta|}{1-\beta}\right), & \beta<q<1,\\[8pt]
0, & \text{иначе}.
\end{cases}
\]

Интерпретация:
\begin{itemize}
  \item при $q<\beta$ значение отрицательное: это зона ближнего отталкивания;
  \item при $\beta<q<1$ сила имеет треугольный профиль с амплитудой $\alpha$;
  \item при $q\ge 1$ взаимодействие обнуляется;
  \item если $\alpha>0$, средняя зона притягивает частицу $i$ к частице $j$;
  \item если $\alpha<0$, средняя зона отталкивает.
\end{itemize}

Парная сила, действующая на частицу $i$ со стороны частицы $j$:
\[
F_{ij}=r_{\max}\,\Phi(q_{ij};\alpha_{ij},\beta)\,e_{ij}.
\]

\section*{6. Изменение цвета}

Парный вклад в изменение цвета частицы $i$:
\[
\Delta c_{ij}=g_{ij}\,\max\left(0,1-\frac{\rho_{ij}}{r_{\max}}\right).
\]

Цвет меняется только внутри радиуса $r_{\max}$ и линейно затухает к нулю на
границе радиуса взаимодействия.

Важная деталь реализации: в суммирование входят также пары $i=j$. Для силы
самовклад фактически равен нулю, потому что $\delta_{ii}=0$ и $e_{ii}=0$.
Для цвета самовклад может быть ненулевым:
\[
\Delta c_{ii}=g_\theta(c_i,c_i).
\]

\section*{7. Агрегирование взаимодействий}

Для каждой частицы:
\[
a_i=\frac{1}{m}\sum_{j=1}^{N}F_{ij},
\]
\[
\dot{c}_i=\sum_{j=1}^{N}\Delta c_{ij}.
\]

В текущей реализации $m$, $\beta$, $h$, $r_{\max}$ и $\Delta t$ являются
скалярными фиксированными параметрами экземпляра \texttt{ParticleLifePlus}, а не
матрицами или векторами по цветам.

\section*{8. Интегрирование движения}

Скорость демпфируется через коэффициент:
\[
\mu=0.5^{\Delta t/h}.
\]

Обновление скорости и положения:
\[
v_i^{n+1}=\mu v_i^n+\Delta t\,a_i^n,
\]
\[
\tilde{x}_i^{n+1}=x_i^n+\Delta t\,v_i^{n+1}.
\]

Это полуявная схема Эйлера: положение обновляется уже новой скоростью.

\section*{9. Граничные условия}

После вычисления $\tilde{x}^{n+1}$ применяется одно из двух правил.

\subsection*{Тор}

Для \texttt{border='torus'}:
\[
x_i^{n+1}=\tilde{x}_i^{n+1}\bmod L,\qquad
v_i^{n+1}\ \text{не изменяется}.
\]

\subsection*{Стенки}

Для \texttt{border='wall'} покомпонентно:
\[
\tilde{x}_{i,k}<0 \Rightarrow
x_{i,k}^{n+1}=-\tilde{x}_{i,k},\quad
v_{i,k}^{n+1}=-v_{i,k}^{n+1},
\]
\[
\tilde{x}_{i,k}>L \Rightarrow
x_{i,k}^{n+1}=2L-\tilde{x}_{i,k},\quad
v_{i,k}^{n+1}=-v_{i,k}^{n+1}.
\]

Затем координата дополнительно обрезается в интервал $[0,L]$.

\section*{10. Интегрирование цвета}

Если \texttt{update\_colors=True}, цвет обновляется как:
\[
\tilde{c}_i^{n+1}=c_i^n+\Delta t\,\dot{c}_i^n,
\]
\[
c_i^{n+1}=
\frac{\tilde{c}_i^{n+1}}{\|\tilde{c}_i^{n+1}\|_2}.
\]

Если \texttt{update\_colors=False}, то:
\[
c_i^{n+1}=c_i^n.
\]

Нормировка удерживает цветовые состояния на единичной сфере
$\mathbb{S}^{K-1}$, если численно не возникает нулевой вектор.

\section*{11. Полный шаг системы}

Один шаг \texttt{step\_state} можно записать как отображение:
\[
S_{n+1}=T_\theta(S_n),
\]
где $T_\theta$ определяется следующей последовательностью:
\begin{enumerate}
  \item для всех упорядоченных пар $(i,j)$ вычислить $\delta_{ij}$, $\rho_{ij}$ и $e_{ij}$;
  \item применить нейросеть к $(c_i,c_j)$ и получить $(\alpha_{ij},g_{ij})$;
  \item вычислить $F_{ij}$ и $\Delta c_{ij}$;
  \item просуммировать силы и цветовые вклады по $j$;
  \item обновить скорость с демпфированием;
  \item обновить положение;
  \item применить граничные условия;
  \item при необходимости обновить и нормировать цвет.
\end{enumerate}

Вызов \texttt{step\_state} принимает \texttt{rng}, но в текущей реализации
случайность в шаге не используется. После инициализации динамика детерминирована
при фиксированных $\theta$ и начальном состоянии.

\section*{12. Рендеринг}

Рендеринг строит RGB-изображение размера \texttt{img\_size}$\times$\texttt{img\_size}.
Фон задается строкой \texttt{background\_color}. Для Particle Life Plus текущая
палитра \texttt{color\_palette} в рендеринге не используется.

Цвет частицы берется из первых трех компонент непрерывного цветового вектора:
\[
\mathrm{rgb}_i=\frac{c_i^{1:3}+1}{2}.
\]

Радиус отрисовки:
\[
R_i=\sqrt{m}\,r_{\mathrm{render}},
\]
где $r_{\mathrm{render}}=\texttt{render\_radius}$.

Для пикселя с координатой $p$ и расстоянием $d_i(p)=\|p-x_i\|_2$ коэффициент
смешивания:
\[
w_i(p)=1-\frac{1}{1+\exp(-s(d_i(p)-R_i))},
\qquad
s=\frac{\texttt{sharpness}}{r_{\mathrm{render}}}.
\]
Эквивалентно, $w_i(p)$ близок к $1$ внутри круга и к $0$ вне круга.

Частицы последовательно накладываются на изображение:
\[
I_{\mathrm{new}}(p)=w_i(p)\,\mathrm{rgb}_i+
(1-w_i(p))\,I_{\mathrm{old}}(p).
\]

\section*{13. Отличие от классического Particle Life}

\begin{longtable}{@{}p{0.30\linewidth}p{0.32\linewidth}p{0.32\linewidth}@{}}
\toprule
Аспект & Particle Life & Particle Life Plus \\
\midrule
\endhead
Цвет & Дискретный индекс $c_i\in\{1,\ldots,K\}$ & Непрерывный нормированный вектор $c_i\in\mathbb{S}^{K-1}$ \\
Параметры силы & Таблицы $\alpha_{c_i,c_j}$, $\beta_{c_i}$, $r_{\max,c_i}$ & Нейросеть $\theta(c_i,c_j)$ возвращает $\alpha_{ij}$ \\
Эволюция цвета & Цвет не меняется & Цвет меняется через $g_\theta(c_i,c_j)$ \\
Инициализация позиций & Может зависеть от распределения по бинам \texttt{x\_dist} & Равномерная по всей области $[0,L]^2$ \\
Масса & Может зависеть от цвета & Один фиксированный скаляр $m$ \\
Граница & Периодическая & Тор или отражающие стенки \\
\bottomrule
\end{longtable}

\section*{14. Вычислительная сложность}

Реализация явно строит все парные взаимодействия через двойной \texttt{vmap}.
Поэтому на один шаг:
\[
\text{время}=O(N^2),\qquad
\text{память}=O(N^2(d+K)).
\]

Пространственная индексация, отсечение соседей до запуска нейросети или
разреженное хранение взаимодействий в текущей версии не используются.

\section*{15. Практические значения по умолчанию}

У самого класса \texttt{ParticleLifePlus} значения по умолчанию:
\[
N=5000,\quad K=6,\quad \Delta t=0.002,\quad h=0.04,\quad
r_{\max}=0.1,\quad m=0.1,\quad \beta=0.3,\quad L=1.
\]

В фабрике ASAL для имени субстрата \texttt{plife\_plus} используются другие
рабочие значения:
\[
N=1000,\quad \Delta t=0.02,\quad
r_{\mathrm{render}}=0.04,\quad
\texttt{sharpness}=30,\quad
\texttt{border='torus'}.
\]

\section*{16. Краткая формула модели}

Particle Life Plus можно сжато определить как нейронно-параметризованную
частичную систему:
\[
\boxed{
\begin{aligned}
\alpha_{ij},g_{ij} &= f_\theta(c_i,c_j),\\
F_{ij} &= r_{\max}\,\Phi(\|x_j-x_i\|/r_{\max};\alpha_{ij},\beta)\,
\frac{x_j-x_i}{\|x_j-x_i\|+\varepsilon},\\
\dot{c}_i &= \sum_j g_{ij}\max(0,1-\|x_j-x_i\|/r_{\max}),\\
v_i^{n+1} &= 0.5^{\Delta t/h}v_i^n+\frac{\Delta t}{m}\sum_j F_{ij},\\
x_i^{n+1} &= \mathrm{boundary}(x_i^n+\Delta t\,v_i^{n+1}),\\
c_i^{n+1} &= \mathrm{normalize}(c_i^n+\Delta t\,\dot{c}_i).
\end{aligned}}
\]

При тороидальной границе во всех формулах расстояние $x_j-x_i$ понимается как
минимальный периодический образ.

\end{document}
